\documentclass[11pt]{article}
\usepackage[margin=1in]{geometry}
\usepackage{hyperref}
\usepackage{xcolor}
\usepackage{enumitem}

\emergencystretch=4em
\tolerance=9999
\hbadness=10000
\sloppy

\title{Metadata Notes for \texttt{gradients\_by\_period/}\\
       \large URI OPeNDAP server, \texttt{sst-aqua.gso.uri.edu}}
\author{An outside reader (grep-dap test case)}
\date{2026-05-17}

\begin{document}
\maketitle

\paragraph{Context.}
This note summarizes the metadata gaps an external client
(\texttt{grep-dap}) found in the public monthly gradient files at
\texttt{https://sst-aqua.gso.uri.edu/opendap/gradients\_by\_period/}
(270 files, \texttt{augmented\_monthly\_stats\_MM\_YYYY.nc},
2003-01 through 2024-12). The headline finding is simple: \textbf{the
files declare no semantic metadata at all.} There are no global
attributes, no per-variable attributes, no coordinate variables, no
fill values. A reader has to infer everything from variable names,
the data themselves, and the companion \texttt{SST\_Orbits/} branch.
Full audit (with inferences classified by confidence): \texttt{docs/uri\_test\_case\_3.tex}.

\section*{What an external reader can infer with high confidence}

These were verified by direct data probes; treat them as facts a
re-issue of the files should encode explicitly.

\begin{enumerate}[leftmargin=*,topsep=2pt]
  \item \textbf{Coordinate grid.}
    Cell-centered $1^\circ \times 1^\circ$ global, lat axis runs
    south-to-north, lon axis runs $-180^\circ$ to $+180^\circ$
    cell-centered. Verified by classifying six known-land points
    (Sahara, Australia interior, Amazon, Central US, Greenland,
    Tibet) and six known-ocean points against
    \texttt{day\_pixel\_count}; only this convention gets all 12
    correct.
  \item \textbf{Variable units.} SST sums in
    $^\circ$C (squared sums in $^\circ$C$^{2}$); gradient sums in
    $^\circ$C\,km$^{-1}$; counts dimensionless. Inherited from the
    sibling \texttt{SST\_Orbits/} branch where the same physical
    quantities are stated.
  \item \textbf{Fill value.} \texttt{NaN}. Observed in every sum
    variable.
  \item \textbf{Temporal coverage.} One calendar month per file,
    encoded in the filename.
  \item \textbf{Day / night split.} Aqua ascending node
    ($\sim$13:30 LST) vs.\ descending node ($\sim$01:30 LST). The
    standard MODIS-Aqua convention; the file does not declare it.
  \item \textbf{Divisor pairing.} Three pixel-count flavours
    paired with three different sum families:
    \begin{itemize}[topsep=0pt]
      \item \texttt{*\_sum\_SST}, \texttt{*\_sum\_eastward\_gradient},
        \texttt{*\_sum\_northward\_gradient},
        \texttt{*\_sum\_magnitude\_gradient} (and their
        \texttt{\_squared}): divide by \texttt{*\_pixel\_count}.
      \item \texttt{*\_sum\_grad\_as\_per\_km} (and
        \texttt{\_squared}): divide by \texttt{*\_as\_pixel\_count}.
      \item \texttt{*\_sum\_grad\_at\_per\_km} (and
        \texttt{\_squared}): divide by \texttt{*\_at\_pixel\_count}.
      \item \texttt{*\_sum\_grad\_mag\_per\_km}: divide by
        \texttt{*\_at\_pixel\_count} (verified empirically: that
        divisor gives a per-cell mean magnitude that matches
        \texttt{*\_sum\_magnitude\_gradient}/\texttt{*\_pixel\_count}).
    \end{itemize}
  \item \textbf{Upstream gradient scale.} The per-pixel gradients
    that get summed here were computed at the L2 swath level with
    an effective operator support of $5\pm 2$ km. Three
    independent estimates (pixel-count ratios, magnitude-RMSE
    matching, autocorrelation length) all converge on this
    number. See \texttt{docs/uri\_test\_case\_4.tex}.
\end{enumerate}

\section*{Essential additions for the file to be self-describing}

\begin{enumerate}[leftmargin=*,topsep=2pt]
  \item \textbf{Coordinate variables.} Add 1-D
    \texttt{lat(180)} and \texttt{lon(360)} with values
    $-89.5,\ldots,+89.5$ and $-179.5,\ldots,+179.5$ and
    attributes
    \begin{quote}\ttfamily
      lat: units="degrees\_north", standard\_name="latitude",
      axis="Y"\\
      lon: units="degrees\_east",  standard\_name="longitude",
      axis="X"
    \end{quote}
    Without these, no CF-aware tool will plot the data without
    user intervention.
  \item \textbf{\texttt{\_FillValue=NaN} on every variable.}
  \item \textbf{Per-variable \texttt{units} and \texttt{long\_name}.}
    Mechanical from the naming grammar; a complete list is in
    \texttt{docs/uri\_test\_case\_3.tex} \S6.
  \item \textbf{Minimum global attribute set.}
    \begin{quote}\ttfamily
      Conventions = "CF-1.8"\\
      title       = "Monthly day/night MODIS Aqua SST and \\
      \phantom{title       = }SST-gradient statistics on a \\
      \phantom{title       = }1$^\circ\times$1$^\circ$ global grid"\\
      institution = (as in SST\_Orbits/)\\
      source      = "MODIS Aqua L2 SST (URI processing); see \\
      \phantom{source      = }SST\_Orbits/"\\
      cdm\_data\_type = "Grid"\\
      time\_coverage\_start / end (per file, from filename)\\
      geospatial\_lat\_min/max = -90 / +90 \\
      geospatial\_lon\_min/max = -180 / +180\\
      geospatial\_lat\_resolution = "1 degree (bin size)"\\
      geospatial\_lon\_resolution = "1 degree (bin size)"
    \end{quote}
  \item \textbf{\texttt{cell\_methods} on every sum variable.}
    For example:
    \begin{quote}\ttfamily
      *\_sum\_SST:\\
      \phantom{xx}cell\_methods = "area: sum time: sum"\\
      *\_sum\_SST\_squared:\\
      \phantom{xx}cell\_methods = "area: sum\_of\_squares time: sum"\\
      *\_pixel\_count:\\
      \phantom{xx}cell\_methods = "area: count time: sum"
    \end{quote}
    These encode the bookkeeping that lets a downstream tool
    recover mean and variance.
  \item \textbf{Comments encoding the inferences in items 5--7 of
    the previous list.} One- or two-line \texttt{comment} fields
    on the affected variables:
    \begin{itemize}[topsep=0pt]
      \item On every \texttt{day\_*} / \texttt{night\_*}:
        ``MODIS Aqua ascending'' / ``descending node, $\sim$13:30
        ($\sim$01:30) local solar time''.
      \item On every \texttt{*\_grad\_as\_per\_km},
        \texttt{*\_grad\_at\_per\_km}, \texttt{*\_grad\_mag\_per\_km}:
        ``along-scan / along-track / swath-relative-magnitude
        component, in the swath-relative frame.''
      \item Global-level comment: \\
        ``Per-pixel gradients were computed at an effective
        spatial scale of $\sim$5 km on L2 swath; values here are
        accumulated into 1$^\circ$ monthly bins.''
      \item Per-variable comment encoding the divisor pairing
        in item~6 above.
    \end{itemize}
\end{enumerate}

\section*{Recommended additions (high payoff for users)}

\begin{enumerate}[leftmargin=*,topsep=2pt,start=7]
  \item \textbf{Scalar \texttt{time} coordinate per file.} Carrying
    a CF-style time avoids forcing every reader to parse the
    filename:
    \begin{quote}\ttfamily
      time = days since 1970-01-01, set to mid-month UTC \\
      time: units="days since 1970-01-01 00:00:00 UTC", \\
      \phantom{time: }standard\_name="time", \\
      \phantom{time: }calendar="gregorian", \\
      \phantom{time: }bounds="time\_bnds"
    \end{quote}
    With \texttt{time\_bnds = [first\_day\_of\_month,
    last\_day\_of\_month]}.

  \item \textbf{ACDD discovery block} (inheritable from
    \texttt{SST\_Orbits/}): \texttt{creator\_name},
    \texttt{creator\_email}, \texttt{creator\_url},
    \texttt{publisher\_name}, \texttt{publisher\_url},
    \texttt{license}, \texttt{acknowledgement},
    \texttt{keywords} / \texttt{keywords\_vocabulary},
    \texttt{platform="Aqua"}, \texttt{instrument="MODIS"},
    \texttt{processing\_level="L4"} (the monthly grid is L4).

  \item \textbf{\texttt{bounds} variables for the coordinate axes:}
    \texttt{lat\_bnds(180, 2)},
    \texttt{lon\_bnds(360, 2)}.

  \item \textbf{\texttt{coordinates} attribute on every data
    variable:}
    \begin{quote}\ttfamily
      coordinates = "time lat lon"
    \end{quote}
\end{enumerate}

\section*{Items not safely guessable from the file}

These an outside reader simply cannot supply; they have to come
from the producing code or notebook:

\begin{itemize}[leftmargin=*,topsep=2pt]
  \item \texttt{date\_created} per file.
  \item \texttt{history} (processing software / version / git
    commit, system, dependencies).
  \item Exact functional form of the upstream gradient operator
    (centered diff vs.\ Sobel vs.\ local least-squares of similar
    support). The effective support is recoverable, but the
    operator shape is not.
  \item Whether \texttt{regridded\_sst} in the orbit files (the
    field on which the gradient is computed) involves a spatial
    averaging step in addition to the high-gradient masking.
    Currently invisible from the public side.
  \item Any minimum-pixel-per-cell suppression rule that may have
    been applied.
\end{itemize}

\paragraph{Bottom line.}
\texttt{gradients\_by\_period/} is a valuable derived product. The
files themselves are unfortunately at the far end of the OPeNDAP
metadata distribution: structurally regular but semantically
silent. The cost of adding the items in the first two lists above
is small (one-time work in the producing script) and the benefit
is that the data become usable by any CF-aware reader without an
out-of-band specification.

\paragraph{Source.}
Full audit and per-variable inferences: \\
\texttt{docs/uri\_test\_case\_3.tex}. \\
Spatial-scale estimation for the underlying gradients: \\
\texttt{docs/uri\_test\_case\_4.tex}.

\end{document}
